\chapter{两节点回路与 81 种符号拓扑}
增加第二个变量后，轨道从浓度轴进入平面。新节点可以只是接收上游信号，也可以将作用传回上游。先比较这两种连接，才能看出闭环符号如何影响扰动，以及二维模型是否已经足以振荡。
\section{信号传递、闭环稳定性与振荡限制}
\begin{equation}\dot x=\alpha_x-\gamma_xx,\qquad\dot y=G(x)-\gamma_yy.\end{equation}
$x$ 的稳态先独立确定，再决定 $y$。纯前向连接不能让 $y$ 反过来影响 $x$；它改变输出幅度、响应时间和滤波性质，而不是构成闭合反馈回路。

这一连接的长期状态可以按上游、下游的顺序确定。若 $G$ 连续非负，则 $x^*=\alpha_x/\gamma_x$、$y^*=G(x^*)/\gamma_y$ 唯一。上游的小偏差按速率 $\gamma_x$ 衰减，下游自身的偏差按速率 $\gamma_y$ 衰减。要进一步确认下游最终确实跟到对应的稳态，可以直接看它的解：
\[
y(t)=e^{-\gamma_yt}y_0+\int_0^t e^{-\gamma_y(t-s)}G(x(s))\dd s.
\]
当 $x(s)\to x^*$ 时，$G(x(s))$ 也趋于 $G(x^*)$。积分中较早的输入影响会随时间衰减，较晚的输入又已接近稳态值，因此 $y(t)\to y^*$。如果上游在稳态附近有小幅、反复的变化，设重复一次所需的时间为 $T$，角频率为 $\omega=2\pi/T$，那么上游变化 $\delta x$ 与下游变化 $\delta y$ 的幅度和时间滞后关系为
$\delta y(i\omega)/\delta x(i\omega)=G'(x^*)/(\gamma_y+i\omega)$：高频幅值被抑制，相位发生滞后。这个关系把上游浓度变化作为输入；若改变的是上游产生率，则还要考虑上游浓度本身需要多久才能跟上。


最简单的无自反馈模型是
\begin{equation}\dot x=f(y)-\gamma_xx,\qquad\dot y=g(x)-\gamma_yy.\label{eq:two}\end{equation}
两条边可以是 $++$（互激活）、$--$（互抑制）、$+-$ 或 $-+$（一激活一抑制）。要知道轻微偏离稳态后会怎样，可以对方程作\emph{线性化}，也就是用稳态附近的斜率近似曲线，只保留小偏差的一阶影响，把它们整理成下面的矩阵 $J$，通常称为雅可比矩阵。对角项表示每个变量自身的恢复作用，另外两项表示一个变量怎样影响另一个：
\begin{equation}J=\begin{pmatrix}-\gamma_x&f'(y^*)\\g'(x^*)&-\gamma_y\end{pmatrix},\quad \operatorname{tr}J=-\gamma_x-\gamma_y<0,\quad\det J=\gamma_x\gamma_y-f'g'.\end{equation}
正反馈时 $f'g'>0$，两个节点会沿回路放大某些偏差。当 $f'g'>\gamma_x\gamma_y$，这股放大作用超过降解的恢复作用，至少有一种偏差会越变越大；另一个方向的偏差仍会恢复，这样的稳态叫“鞍点”。是否还存在另外两个稳定状态，需要继续寻找所有产生与降解平衡的位置。负反馈时 $f'g'<0$，这个模型的稳态附近偏差会衰减，有时还会边衰减边上下摆动。

这里判断的是偏差随时间放大还是缩小。偏差可以分成两种基本变化方式，每一种都带有 $e^{\lambda t}$ 这样的时间因子，$\lambda$ 称为特征值。$\lambda<0$ 表示这一种偏差衰减，$\lambda>0$ 表示它增长；若 $\lambda$ 是复数，偏差还会在增长或衰减的同时摆动，增长还是衰减由它的实部决定。令 $q=f'(y^*)g'(x^*)$，把两种变化方式的特征值写出来，得到
\[
\lambda_{\pm}=-\frac{\gamma_x+\gamma_y}{2}
\pm\sqrt{\frac{(\gamma_x-\gamma_y)^2}{4}+q}.
\]
这两个数满足 $\lambda^2+(\gamma_x+\gamma_y)\lambda+\gamma_x\gamma_y-q=0$。前面写出的 $\operatorname{tr}J$ 称为\emph{迹}，是对角项之和；$\det J$ 称为\emph{行列式}，对这里的二维矩阵就是两个对角项的乘积减去两个交叉项的乘积。它们分别对应两个特征值的和与积，因此能用来判断偏差会怎样变化。若 $\det J>0$，两个实数特征值的和为负、积为正，所以都为负；若它们是复数，共同的实部为 $-(\gamma_x+\gamma_y)/2$，也为负。两种情况都表示偏差会衰减。若 $\det J<0$，两个特征值一正一负，出现鞍点。若 $\det J=0$，至少有一个特征值为零，仅看稳态附近的斜率还不能判断这个方向会怎样。
负反馈的衰减振荡条件进一步是
$q<-(\gamma_x-\gamma_y)^2/4$，所以负反馈符号本身甚至不保证\emph{衰减}振荡；正反馈 $q>0$ 则两种变化方式都不会带有这种摆动。
\begin{figure}[H]\centering
\begin{tikzpicture}
\foreach \i/\title in {0/互激活,1/互抑制,2/负反馈}{
\node at (4.1*\i+1,1.2) {\title};
\node[gene] (x\i) at (4.1*\i,0) {$X$};
\node[gene] (y\i) at (4.1*\i+2,0) {$Y$};}
\draw[act] (x0) to[bend left=25] (y0);\draw[act] (y0) to[bend left=25] (x0);
\draw[rep] (x1) to[bend left=25] (y1);\draw[rep] (y1) to[bend left=25] (x1);
\draw[act] (x2) to[bend left=25] (y2);\draw[rep] (y2) to[bend left=25] (x2);
\end{tikzpicture}
\caption{两个交叉边形成的反馈。互激活与互抑制均为正回路；一激活一抑制为负回路。交换节点得到另一种带标签的负回路。}
\end{figure}

在负反馈下，局部稳定还可以与稳态的唯一性联系起来。假设 $f,g$ 非负、有界且为 $C^1$，并且一个随输入增加而严格上升，另一个严格下降。在稳态方程中消去 $y$，得到
\[
x=T(x):=\frac1{\gamma_x}f\!\left(\frac{g(x)}{\gamma_y}\right).
\]
经过一正一负两次调控，$T$ 不会随 $x$ 增加而上升，而且有上限。所以 $T(x)-x$ 一直下降，在零处非负，$x$ 足够大时为负，恰好经过零一次。随后 $y=g(x)/\gamma_y$ 也唯一。又有 $q\leq0$，由上面的计算，两种偏差都只能衰减，因此这个唯一稳态在小偏差后能逐渐恢复。这与单节点自抑制相呼应：闭环总体为负时，偏差受到恢复作用；两个变量则额外允许恢复过程带有衰减振荡。


边衰减边摆动，仍然是逐渐回到稳态；持续振荡则要求轨迹每隔一段时间回到原处，并重复同样的变化。对式~\eqref{eq:two}，可以排除后一种情况，原因是这个模型始终使浓度平面中的面积收缩。这是\emph{Bendixson 判据}在此模型中的应用：若一个没有空洞的平面区域内，散度始终严格为同一符号，就不存在非恒定周期轨道。下面解释为什么面积收缩会产生这个限制。

把两个变化率写成 $\bm v=(f(y)-\gamma_xx,g(x)-\gamma_yy)$，则沿各自浓度方向计算斜率并相加，有
\[
\frac{\partial\dot x}{\partial x}+\frac{\partial\dot y}{\partial y}=-\gamma_x-\gamma_y<0.
\]
这个和称为“散度”。它为负，表示一小片初始状态经过同样一段时间的演化后，占据的面积会缩小。

假如存在一条不断重复的轨迹，它在浓度平面上就围成闭合曲线 $C$，内部区域记为 $\Omega$。轨迹不会自相交，因为同一组浓度只能对应一种后续变化。还假设这条曲线及其内部都处于方程可以正常使用的范围内。数学上常通过\emph{单连通区域}保证这一点：直观地说，区域中没有空洞，区域内的闭合曲线所围住的内部也属于这个区域。沿曲线，变化方向与曲线相切，没有穿过边界向外的净流动；但内部处处收缩又要求净流动小于零。平面散度定理把“沿边界的净流动”与“内部收缩率的总和”联系起来，因此这两种说法可以写成同一个面积关系：
\[
0=\oint_C\bm v\cdot\bm n\,\dd\ell
=\iint_\Omega\nabla\cdot\bm v\,\dd x\dd y
=-(\gamma_x+\gamma_y)\operatorname{Area}(\Omega)<0,
\]
相互矛盾。左边沿边界计算净向外流动，右边把内部的收缩率加起来；$\bm n$ 表示指向边界外侧的方向。因此该模型没有持续重复的非恒定轨迹。

这里用到了“没有自调控”和“一阶降解”。换成其他二维负反馈模型，收缩率未必始终为负，结论也可能改变。另外，没有持续振荡并不等于所有初始状态都回到同一个稳态：正反馈仍可能形成多个长期状态。

\section{81 种拓扑：从结构筛选到动力学机制}
前面的模型只保留两条交叉边。允许两个节点分别自激活、自抑制或没有自边后，产生率变为 $F_x(x,y),F_y(x,y)$，一般形式是
\begin{equation}
\dot x=F_x(x,y)-\gamma_xx,\qquad
\dot y=F_y(x,y)-\gamma_yy.\label{eq:two_general}
\end{equation}
以下分类假设产生率非负、有界且为 $C^1$，边的符号在所讨论的浓度范围内固定；缺失边表示相应偏导数恒为零。输入固定，没有延迟，也没有其他隐藏动态变量。这样的约定让“一个拓扑”具有可比较的含义；若边会随浓度改变符号，就不属于同一个固定符号类别。

四条边 $(X\to X,X\to Y,Y\to X,Y\to Y)$ 各有 $0,+,-$ 三种选择，所以带标签的情况恰为 $3^4=81$。交换 $X,Y$ 后不变的图需要两条自边相同、两条交叉边相同，共 $3^2=9$ 个；其余图两两配对，故有 $9+(81-9)/2=45$ 个无标签类别。本文先保留标签，方便区分输入、输出及两个节点的不同时间尺度。

逐一背 81 张图不利于理解。更有效的办法是先看两个节点是否互相连通，再看有没有自激活：没有交叉边有 $9$ 种；恰有一条交叉边有 $2\times2\times9=36$ 种；两条交叉边都存在时，同号与异号各有 $2\times9=18$ 种。每种连接下，自边的 $9$ 种组合中有 $4$ 种完全没有正自边、$5$ 种至少有一条正自边。因此所有情况恰好分成下表的八组。

\begin{center}\small
\begin{tabular}{c r >{\raggedright\arraybackslash}p{4.6cm} >{\raggedright\arraybackslash}p{7.6cm}}
\toprule 组 & 数目 & 结构 & 先分析什么\\\midrule
R1 & 4 & 无交叉边，无自激活 & 两个独立的恢复过程；每个节点有唯一稳态。\\
R2 & 5 & 无交叉边，有自激活 & 独立记忆单元；自激活足够强时可有多个表达组合。\\
R3 & 16 & 单向连接，无自激活 & 有限增益的信号传递；唯一稳态，不能自行持续振荡。\\
R4 & 20 & 单向连接，有自激活 & 上游状态传给下游，或下游自己存储状态；先后求解两个节点。\\
R5 & 8 & 同号交叉边，无自激活 & 互激活或互抑制开关；多稳态取决于闭环增益。\\
R6 & 10 & 同号交叉边，有自激活 & 多个正反馈共同作用；考察状态数、阈值与迟滞，不能经 Hopf 失稳。\\
R7 & 8 & 异号交叉边，无自激活 & 唯一局部渐近稳定的稳态，可有衰减摆动，无周期轨道。\\
R8 & 10 & 异号交叉边，有自激活 & 正放大与负恢复竞争；可能出现 Hopf 分岔，也可能只是恢复或多稳态。\\
\bottomrule
\end{tabular}\end{center}

这些是\emph{结构分组与解析筛选}，不是把每张图标成固定功能。R1 的恢复看起来简单，却可以提供稳定的基线；R3 可以滤去快速噪声；R2、R4 的记忆也不必依赖跨节点闭环。哪些情况“有趣”，取决于想解释的是恢复、状态选择、时间响应还是振荡。图~\ref{fig:81map}把 81 个组合完整放进同一张索引图；附录还给出逐张回路图与编号，方便查到任何一种情况。
\begin{figure}[H]\centering
\includegraphics[width=.76\linewidth]{figures/topology_classes.pdf}
\caption{完整的 81 种两节点符号拓扑。行是 $(X\to X,Y\to Y)$ 的自边组合，列是 $(X\to Y,Y\to X)$ 的交叉边组合；$0$ 表示缺失。每个格子的 R 编号对应正文八组，颜色表示分组，不表示已经实现某种功能。}\label{fig:81map}
\end{figure}


R1 与 R2 没有节点间连接，整个系统是两个一维系统并列。R1 的自边只有缺失或抑制，前面的一维恢复论证分别适用。R2 则允许其中一个或两个节点具有自激活；例如两个节点都在双稳态参数区时，每个节点各有低、高两种稳定状态，组合后可以有四种稳定表达状态。这是两份独立记忆，并不需要互抑制。

R3 与 R4 虽然有传递，但作用只有一个方向。先求不受下游影响的节点，再把它的稳态代入下游方程。R3 的两条自边都不会促进自身产生，因此两个步骤各有唯一稳态；上游趋于稳态后，下游也会逐渐跟上。严格地看，下游偏差的恢复率至少为 $\gamma_y$，而上游尚未稳定造成的附加输入会趋于零，所以前面的积分估计仍保证它最终恢复。R4 的任一步都可能因自激活而出现多个解：需要分别问“上游能否选择状态”和“同一上游输入下，下游能否保留不同状态”。这两种记忆机制值得区分。

这四组总共有 $45$ 种，都没有跨节点闭环。它们不能在固定输入下产生非恒定周期轨道：若整个系统周期变化，上游也必须周期变化；一维自治方程迫使上游恒定，随后下游又成为一个一维自治方程，也只能恒定。因此“有两个变量”本身不是振荡机制。

对 R5 与 R6，同号交叉边构成正回路。互激活增强共同表达的变化，互抑制则增强两者之间的差异；它们与自激活一样，都能为状态选择提供放大作用，但还需要足够陡的产生率曲线。一个有用的结构统计是：含至少一条正自边的图有 $81-2^2\times3^2=45$ 种；没有正自边却有正交叉回路的图有 $2\times4=8$ 种，所以含正回路的共有 $53$ 种。

其余 $28$ 种就是 R1、R3、R7。前两组已知只有一个稳态，下面说明 R7 也如此。因此在本文的固定符号、有限产生率模型中，多稳态只能在这 $53$ 种含正回路的图里寻找。这个条件只是筛选入口，不代表这 $53$ 种在任意参数下都会多稳态。

R7 可以保留自抑制。给定 $y$ 时，$F_x(x,y)-\gamma_xx$ 随 $x$ 严格下降，恰有一个零点 $x=h(y)$；同理 $y=k(x)$ 唯一。由定态方程求导，
\[
h'(y)=\frac{\partial_yF_x}{\gamma_x-\partial_xF_x},\qquad
k'(x)=\frac{\partial_xF_y}{\gamma_y-\partial_yF_y}.
\]
分母均为正，两条交叉边异号，所以 $h(k(x))$ 非增，与直线 $x$ 只相交一次。边界处的解也可以由同样的单调性处理。因而 R7 有唯一稳态；两个对角恢复率为负、交叉作用乘积为负，保证它局部渐近稳定。散度在整个区域为负，又排除了周期轨道。可见无反馈、负反馈和弱正反馈都可能表现为恢复，但原因并不相同。


在稳态附近，式~\eqref{eq:two_general}的雅可比可以写成
\[
J=\begin{pmatrix}a&b\\c&d\end{pmatrix},\quad
 a=\partial_xF_x-\gamma_x,\quad d=\partial_yF_y-\gamma_y,
\quad b=\partial_yF_x,\quad c=\partial_xF_y.
\]
这里 $a,d$ 是自身产生与降解合起来的\emph{净作用}，不是自边符号本身。正自边也可能太弱，以至于 $a$ 或 $d$ 仍为负。两个特征值满足
\[
\lambda_{\pm}=\frac{a+d}{2}\pm\frac12\sqrt{(a-d)^2+4bc}.
\]
这一个公式便能检验全部八组的 Hopf 入口。R1--R4 至少一条交叉边缺失，$bc=0$，两根为实数；R5、R6 有 $bc>0$，根号内严格为正，两根也为实数，所以不会通过一对复特征值穿越虚轴产生 Hopf 分岔。R7 虽然 $bc<0$，但 $a,d<0$，迹不可能到零；没有正自边的所有 $36$ 种图还都满足全域面积收缩。

于是只有 R8 的 $10$ 种具备这种失稳的符号条件：两条交叉边异号，而且至少有一条正自边。在候选临界点，
\[
a+d=0,\qquad ad-bc>0.
\]
代入 $d=-a$，条件变成 $-bc>a^2$。直观地说，一边的自放大刚好抵消另一边的净恢复，而负闭环还必须足够强，才能使偏差以摆动方式变化。这仍只是必要的线性条件，真正的 Hopf 分岔还需参数横截和非线性非退化条件。

为了看到“同一结构可以恢复，也可以振荡”，考虑 R8 中的一张图：$X$ 自激活，$X$ 激活 $Y$，$Y$ 抑制 $X$，$Y$ 没有自边。取
\begin{align}
H(x)&=\frac{x^4}{1+x^4},\\
\dot x&=\frac{a_0+b_0H(x)}{1+y^4}-x,\qquad
\dot y=\varepsilon\{d_0H(x)-y\}.
\end{align}
自激活让 $X$ 快速上升，较慢积累的 $Y$ 随后压低 $X$；$X$ 降低后，$Y$ 的产生也减少，抑制逐渐解除。这就是正放大与负恢复交替工作的机制。固定 $a_0=0.4,b_0=4,d_0=2$，图~\ref{fig:twoosc}比较 $\varepsilon=1$ 与 $\varepsilon=0.05$：前者在小扰动后恢复，后者数值上持续振荡。改变的是 $Y$ 的时间尺度，而不是连边。脚本检查了示例定态的局部稳定性及后期轨迹，图中振荡仍是数值例子，没有替代周期轨道的严格证明。
\begin{figure}[H]\centering\includegraphics[width=.95\linewidth]{figures/two_node_oscillation.pdf}
\caption{同一 R8 拓扑的两种参数行为。左：下游较快，偏差衰减；右：下游较慢，自激活与负反馈轮流主导，数值轨迹持续摆动。仅展示后期轨迹，未标注不必要的数值根。}\label{fig:twoosc}
\end{figure}

二维部分的结论可以这样使用：先凭结构排除不适合的机制，再用反馈增益与时间尺度检验剩余候选。开关重点看交点与正回路；Hopf 振荡重点看 R8、迹与行列式。增加第三个节点，则可以在没有正自边的负反馈环中提供新的动态环节，也可以让输入沿多条路径抵达输出。
